\section*{Limitations}

\paragraph{Readout scope.}
SRP targets final-unembedding alignment: it identifies which sparse readout
directions support or oppose a chosen scalar once the hidden state reaches the
LM head. The feature-resolved DLA analysis composes SRP with additive
residual-stream attribution, providing candidates for intervention studies that
test component or feature necessity.

\paragraph{Dictionary non-uniqueness.}
The sparse basis is not unique: dictionary features can split, merge, rotate, or
depend on recipe choices. We therefore treat reconstruction error, sign
agreement, replacement fidelity, and stability checks as part of the result
rather than as auxiliary diagnostics. Feature names are row-level summaries from
high-activating unembedding rows; the measured objects are feature ids, signed
terms, and residual errors.

\paragraph{Target and model coverage.}
SRP currently handles linear scalar targets at the final readout. Distributional
claims, multi-token answer quality, and sequence-level behavior require chosen
linear proxies or additional evaluation. The strongest qualitative evidence in
the paper comes from Qwen settings with good target fidelity; other families and
diffuse targets can preserve signs while having larger errors. Small exact
margins are especially fragile, so displayed accounts should be read together
with their local reconstruction error.

\paragraph{Tokenizer and benchmark effects.}
Because SRP factorizes vocabulary rows, labels and feature groupings can reflect
tokenizer geometry, token frequency, row norms, special tokens, and surface-form
artifacts. Contrasts and group targets reduce dependence on any one row but do
not remove it. Benchmark-format rows measure whether task-shaped readout targets
are reconstructed, not whether the model solved the underlying task.

\paragraph{Readout edits.}
The readout-edit experiment tests constrained movement of specified lexical
scores. It should not be interpreted as broad behavioral control, safety editing,
or evidence that the same directions have the same causal role throughout the
network. Larger behavioral effects require intervention studies that track
distribution shift, downstream generations, and off-target changes.
